\subsection{形式幂级数环的自同构预备知识}

引理 2. 设 \(f(\mathbf{X}_1, \mathbf{X}_2, \ldots, \mathbf{X}_n)\) 是一个系数属于环 E 的非零形式幂级数。存在整数 \(\neq 0\)\(u(i) \geqslant 1\)（\(1 \leqslant i \leqslant n - 1\)），使得

\[
f (\mathrm{T} ^ {u (1)}, \dots , \mathrm{T} ^ {u (n - 1)}, \mathrm{T}) \neq 0.
\]

假设已经确定整数 \(u(i) \geqslant 1\)（\(1 \leqslant i \leqslant k - 1\)），使得 \(f(\mathbf{X}_n^{u(1)}, \ldots, \mathbf{X}_n^{u(k-1)}, \mathbf{X}_k, \ldots, \mathbf{X}_n) \neq 0\)。我们将确定一个整数 \(u(k) \geqslant 1\)，使得

\[
f \left(\mathrm{X} _ {n} ^ {u (1)}, \dots , \mathrm{X} _ {n} ^ {u (k - 1)}, \mathrm{X} _ {n} ^ {u (k)}, \mathrm{X} _ {k + 1}, \dots , \mathrm{X} _ {n}\right) \neq 0.
\]

于是引理可由归纳法证明。

注意，级数 \(f(\mathbf{X}_n^{u(1)}, \ldots, \mathbf{X}_n^{u(k-1)}, \mathbf{X}_k, \ldots, \mathbf{X}_n)\) 可以看作关于 \(\mathbf{X}_k\) 和 \(\mathbf{X}_n\) 的级数，其系数属于 \(\mathbf{E}[[\mathbf{X}_{k+1}, \ldots, \mathbf{X}_{n-1}]]\)。因此只需在 \(n = 2\) 时证明引理。

因此令

\[
f = \sum_ {i, j} e _ {i j} \mathbf {X} ^ {i} \mathbf {Y} ^ {j} \in \mathrm{E} [ [ \mathbf {X}, \mathbf {Y} ] ]
\]

其中 \(f \neq 0\)。令 \(G \subset N \times N\) 为满足  的有序对 \((i, j)\) 组成的非空集。赋予 \(e_{ij} \neq 0\)\(N \times N\) 字典序。令 \((c, d)\) 为 \(G\) 的最小元。取整数 \(p > d\)。在下式的展开式中

\[
f (\mathbf {T} ^ {p}, \mathbf {T}) = \sum_ {(i, j) \in \mathbb {G}} e _ {i j} \mathbf {T} ^ {i p + j}
\]

我们考察次数为 \(cp + d\) 的项。如果 \(ip + j = cp + d\)，则不可能有 \(i \geqslant c + 1\)，因为这将导致

\[
i p + j \geqslant (c + 1) p + j \geqslant (c + 1) p > c p + d;
\]

同样也不可能有 \(i < c\)，因为 \((c, d)\) 是 \(G\) 的最小元；因此 \(i = c\)，进而 \(j = d\)。所以  中次数为 \(cp + d\) 的项是 \(f(T^p, T)\)\(e_{cd}T^{cp + d}\)。由于 \(e_{cd} \neq 0, f(T^p, T) \neq 0\)，。引理得证。

在环 \(\operatorname{E}[[\mathbf{X}_1, \ldots, \mathbf{X}_n]]\) 中，令 \(a\) 为不含常数项的形式幂级数所成的理想。若 \(w_1, \ldots, w_n\) 是 \(a\) 的元素，记得映射 \(f(\mathbf{X}_1, \ldots, \mathbf{X}_n) \mapsto f(w_1, \ldots, w_n)\) 是环  的唯一自同态 \(s\)，满足 \(\operatorname{E}[[\mathbf{X}_1, \ldots, \mathbf{X}_n]]\)\(s(\mathbf{X}_i) = w_i\)（\(1 \leqslant i \leqslant n\)）（第三章，§ 4，第 5 节，命题 6）。

取 \(w_1 = \mathbf{X}_1 + \mathbf{X}_n^{u(1)}, \ldots, w_{n-1} = \mathbf{X}_{n-1} + \mathbf{X}_n^{u(n-1)}, w_n = \mathbf{X}_n\)，其中 \(u(i)\) 是满足 \(\geqslant 1\) 的整数。令 \(s'\) 为 \(\operatorname{E}[[\mathbf{X}_1, \ldots, \mathbf{X}_n]]\) 的自同态，它把 \(\mathbf{X}_1\) 映为 \(\mathbf{X}_1 - \mathbf{X}_n^{u(1)}, \ldots, \mathbf{X}_{n-1}\)，把  映为 \(\mathbf{X}_{n-1} - \mathbf{X}_n^{u(n-1)}\)，并把 \(\mathbf{X}_n\) 映为 \(\mathbf{X}_n\)。于是对 ，有 \(s'(s(\mathbf{X}_i)) = \mathbf{X}_i\)，所以 \(1 \leqslant i \leqslant n\)\(s' \circ s\) 是恒等自同构；对于 \(s \circ s'\) 也同样如此。因此 \(s\) 是自同构。

引理 3. 设 \(f\) 是 \(\operatorname{E}[[\mathbf{X}_1, \ldots, \mathbf{X}_n]]\) 的非零元。存在整数 \(u(i) \geqslant 1 (1 \leqslant i \leqslant n - 1)\)，使得由下式定义的  的自同构 \(s\)\(\operatorname{E}\)

\[
s (\mathbf {X} _ {i}) = \mathbf {X} _ {i} + \mathbf {X} _ {n} ^ {u (i)} \quad (1 \leqslant i \leqslant n - 1)
\]

并且 \(s(\mathbf{X}_n) = \mathbf{X}_n\) 将 \(f\) 映为一个元 \(g\)，使得 \(g(0, \ldots, 0, \mathbf{X}_n) \neq 0\)。

\(g(0,\ldots,0,\mathbf{X}_n) = f(\mathbf{X}_n^{u(1)},\ldots ,\mathbf{X}_n^{u(n - 1)},\mathbf{X}_n)\)。因此引理 3 是引理 2 的推论。

